\documentclass[11pt,a4paper]{article}

% ---------------------------------------------------------------------------
% Classificacao de topicos - case VDE - Fase 2
% Fonte oficial do relatorio. Compila com pdfLaTeX (duas passadas).
% Estrutura segue o PDF do case: a) sistema, b) os 100 rotulos, c) diagnostico,
% d) producao, e) memo para o professor.
% ---------------------------------------------------------------------------

\usepackage[T1]{fontenc}
\usepackage[utf8]{inputenc}
\usepackage[brazilian]{babel}
\usepackage[margin=1.9cm]{geometry}
\usepackage{booktabs}
\usepackage[table]{xcolor}
\usepackage{microtype}
\usepackage{enumitem}
\usepackage{titlesec}
\usepackage{tcolorbox}
\tcbuselibrary{skins}
\usepackage{tikz}
\usetikzlibrary{arrows.meta,positioning}
\usepackage{fancyhdr}
\usepackage{lastpage}
\usepackage{needspace}
\usepackage[scaled=1.0]{helvet}
\renewcommand{\familydefault}{\sfdefault}
\usepackage[hidelinks]{hyperref}

\definecolor{azul}{HTML}{1F4E79}
\definecolor{cinza}{HTML}{F2F4F7}

\newcommand{\bd}[1]{\textbf{#1}}
\setlist[itemize]{leftmargin=1.2em,itemsep=1pt,topsep=2pt}
\setlist[enumerate]{leftmargin=1.4em,itemsep=1pt,topsep=2pt}
\setlength{\parskip}{3.5pt}
\setlength{\parindent}{0pt}

\titleformat{\section}{\normalfont\large\bfseries\color{azul}}{}{0pt}{}[{\vspace{-0.15em}\color{azul}\titlerule[0.8pt]}]
\titlespacing*{\section}{0pt}{8pt}{3pt}

\fancypagestyle{rodape}{%
  \fancyhf{}%
  \renewcommand{\headrulewidth}{0pt}%
  \fancyfoot[L]{\small\color{black!55}Classificação de tópicos: case VDE (Fase 2)}%
  \fancyfoot[R]{\small\color{black!55}Página \thepage\ de \pageref{LastPage}}%
}
\pagestyle{rodape}

\newcommand{\kpi}[2]{%
  \fcolorbox{azul!55}{white}{\parbox[c][12mm][c]{0.215\textwidth}{\centering
    {\Large\bfseries\color{azul}#1}\\[2pt]{\scriptsize\bfseries\color{black!65}#2}}}%
}

\tikzset{
  box/.style={draw=azul,rounded corners,fill=white,align=center,inner sep=3pt,
              minimum height=8mm,text width=20mm,font=\small},
  arr/.style={-{Latex[length=2mm]},thick,draw=azul}
}

\begin{document}
{\LARGE\bfseries\color{azul}Classificação de tópicos: case VDE\par}
\vspace{2pt}
{\large Fase 2: um sistema que decide sem modelo de linguagem\par}
\vspace{3pt}
{\small Leo Vitor\par}
\vspace{3pt}
{\color{azul}\hrule height 0.8pt}
\vspace{6pt}
\thispagestyle{rodape}

\begin{tcolorbox}[colback=cinza,colframe=azul!35,
                  borderline west={2.5pt}{0pt}{azul},
                  title={\bfseries Sumário executivo},coltitle=azul,colbacktitle=cinza,
                  boxrule=0.6pt,arc=2pt,left=6pt]
\begin{itemize}
\item \bd{Resultado:} publiquei 115 das 659 questões (17,4\%) e mandei 544 para a revisão do
      professor. Nas 559 questões fora da folha-verdade, publiquei 94 (0,17).
\item \bd{Previsão de precisão: 0,76, com faixa de 0,53 a 0,90.} É a medida mais honesta que
      tenho: 13 acertos em 17 publicações, numa avaliação em que a regra nunca viu as questões
      usadas para medi-la. Quando meço a regra nas mesmas questões que a ajustaram, ela acerta
      19 de 20 (0,95), mas esse número é otimista e não é a minha previsão.
\item \bd{A meta de 95\% não é demonstrável com 100 rótulos.} Mesmo sem nenhum erro, eu
      precisaria conferir umas 60 publicações para provar 95\%; na folha-verdade o sistema
      publica 20.
\item \bd{Nenhum modelo de linguagem na hora de decidir.} Um comando refaz toda a entrega em
      cerca de um minuto, sem internet. O único material produzido por um LLM é a descrição
      ampliada dos temas, gerada uma vez e guardada no repositório.
\end{itemize}
\end{tcolorbox}

\vspace{2pt}
\noindent\kpi{17,4\%}{COBERTURA (115/659)}\hfill\kpi{0,76}{PRECISÃO PREVISTA}\hfill
\kpi{0,74}{ASSUNTO CERTO EM 1º}\hfill\kpi{R\$ 0}{CUSTO DE API}

% ================================= a) O SISTEMA ============================
\section{a) O sistema}

O sistema responde a três perguntas, sempre na mesma ordem: qual é o tema mais provável, a que
assunto geral ele pertence e se vale publicar. Chamo de \bd{ramo} o tema de primeiro nível da
árvore (por exemplo, ``Teoria da Pena''); os subtemas ficam abaixo dele. Nenhuma das três perguntas
consulta um modelo de linguagem.

\begin{center}
\begin{tikzpicture}[node distance=5mm]
\node[box] (a) {Limpeza\\\scriptsize HTML $\to$ texto};
\node[box,right=of a] (b) {Busca\\\scriptsize 4 métodos};
\node[box,right=of b] (c) {Regra do pai\\\scriptsize subtema $\to$ ramo};
\node[box,right=of c] (d) {Confiança\\\scriptsize probabilidade};
\node[box,right=of d] (e) {Decisão\\\scriptsize publicar/revisar};
\node[box,above=4mm of b,fill=cinza] (t) {Temas\\\scriptsize enriquecidos};
\draw[arr] (a)--(b); \draw[arr] (b)--(c); \draw[arr] (c)--(d); \draw[arr] (d)--(e);
\draw[arr] (t)--(b);
\end{tikzpicture}
\end{center}

\begin{enumerate}
\item \bd{Temas enriquecidos.} O catálogo traz, para cada tema, só o nome e o caminho na árvore,
      o que é pouco para comparar com o texto de uma questão. Um LLM escreveu, uma única vez, uma
      descrição ampliada dos 351 temas ativos: sinônimos, institutos jurídicos e uma definição
      curta. Um programa conferiu automaticamente campos, leis citadas e repetições entre pai e
      filho, e o arquivo ficou congelado no repositório. Foi a peça que mais rendeu (tabela
      abaixo), e é exatamente o uso de LLM que o case permite.
\item \bd{Busca com quatro métodos.} Cada método ordena os temas do mais ao menos parecido com
      a questão. Dois contam palavras em comum, valorizando as raras (TF-IDF por palavra e o BM25,
      usado em buscadores); um compara pedaços de palavras, o que tolera plural e conjugação
      (TF-IDF por caractere); o quarto junta os três rankings num só (fusão RRF). A fusão escolhe
      o tema; os quatro juntos dizem se há \bd{consenso}.
\item \bd{Regra do pai.} Se o tema mais provável é um subtema, publico o ramo acima dele. Pela
      régua do case, publicar no pai conta como acerto para o subtema certo e para todos os seus
      irmãos, então essa regra nunca transforma um acerto em erro. Na folha-verdade, ela eliminou
      os erros entre subtemas vizinhos.
\item \bd{Confiança.} Para cada questão, calculo 13 \bd{sinais}: números que ajudam a prever se a
      busca acertou. Por exemplo, a distância entre o primeiro e o segundo tema do ranking, quantos
      métodos concordam, o nível do tema na árvore e se a disciplina de origem bate com a do tema.
      Uma regressão logística, um modelo estatístico simples, combina os sinais numa probabilidade
      de acerto. Ela aprendeu com os 100 rótulos, sempre medida em questões que não viu.
\item \bd{Decisão.} Publico só quando pelo menos três dos quatro métodos concordam e a
      probabilidade de acerto passa de 0,896 (o \bd{limiar}). O resto vai para revisão, com a
      sugestão do sistema.
\end{enumerate}

\begin{center}
\small
\begin{tabular}{lcc}
\toprule
Descrição usada para cada tema (92 questões com tema) & Ramo certo em 1º & Diferença \\
\midrule
Caminho e título (o catálogo original) & 0,630 & \\
Catálogo e títulos dos subtemas & 0,630 & 0,000 \\
Catálogo e termos de um dicionário jurídico & 0,620 & $-$0,011 \\
Só sinônimos & 0,609 & $-$0,022 \\
\bd{Descrição enriquecida (sinônimos, institutos e definição)} & \bd{0,739} & \bd{+0,109} \\
\bottomrule
\end{tabular}
\end{center}
\vspace{-4pt}
{\small Escrevi o critério antes do teste: só entraria com ganho de pelo menos 0,05. O ganho foi
de 0,109, com faixa provável de 0,033 a 0,196 comparando questão a questão: positivo, mas sem
folga. Nenhum campo sozinho resolve; o ganho vem da combinação.}

\bd{O que testei e descartei.} Cada ideia só entrava se melhorasse o acerto do ramo em primeiro.
\begin{itemize}
\item \bd{Embeddings}, que representam o texto como uma lista de números que captura o sentido
      (testei o multilingual-e5-large, inclusive com os temas enriquecidos): 0,53 a 0,68 sozinhos,
      e a melhor combinação com a busca ganhou uma única questão. O que separa temas jurídicos
      vizinhos é a palavra exata (um número de lei, um termo técnico), e o embedding aproxima
      justamente o que eu preciso separar.
\item \bd{Vizinhos entre questões (kNN)}, que dá a uma questão o tema das questões rotuladas mais
      parecidas: só 41 dos 109 ramos aparecem na folha-verdade, então nem um acerto perfeito
      passaria de 0,79. Ficou em 0,43 a 0,48.
\item \bd{Um segundo modelo para reordenar os candidatos} (16 sinais por par questão e tema):
      empatou ou perdeu. \bd{Decidir primeiro a disciplina e depois o tema}: deu a mesma decisão
      da regra do pai. \bd{Dar peso à disciplina de origem}: duas questões a mais em 92, dentro do
      acaso.
\item \bd{Um sinal de lei citada}, que entrou e saiu. Tinha um único caso positivo nas 100 e
      dominava a probabilidade. Removê-lo custou 5 das 6 publicações de Legislação Especial e
      derrubou a precisão honesta de 0,87 para 0,76. Preferi um número pior a um sistema apoiado
      num caso só.
\end{itemize}

% ================================= b) OS 100 ROTULOS =======================
\section{b) Os 100 rótulos}

\bd{Como usei.} A folha-verdade (as 100 questões rotuladas pela equipe) teve três papéis, sempre
separados: escolher as peças do sistema, com o critério escrito antes; treinar o modelo de
confiança; e medir a precisão. O cuidado central foi não medir nas mesmas questões em que ajustei.
Medir assim é como corrigir a própria prova: o resultado sai bom demais. Os estatísticos chamam
isso de medida \bd{in-sample} (``dentro da amostra'').

Para evitar isso, usei \bd{validação cruzada aninhada}. Divido as 100 questões em 5 partes; ajusto
o modelo e o limiar em 4 delas e meço na quinta, que ficou de fora; repito trocando a parte de
fora e, depois, refaço tudo com 5 sorteios diferentes. Questões quase iguais ficariam sempre na
mesma parte (não houve nenhuma). Ao auditar esse processo, achei um \bd{vazamento}: o limiar era
escolhido com probabilidades que já tinham visto a parte de teste. Corrigido, as publicações nessa
medida caíram de 33 para 17. O número que entrego é o corrigido.

\begin{center}
\small
\begin{tabular}{lccc}
\toprule
Como a precisão foi medida (régua do case) & Publicadas & Acertos & Precisão (faixa de 95\%) \\
\midrule
Nas mesmas 100 que ajustaram a regra (in-sample) & 20 & 19 & 0,95 (0,76 a 0,99) \\
\bd{Aninhada: regra ajustada sem ver o grupo medido} & \bd{17} & \bd{13} & \bd{0,76 (0,53 a 0,90)} \\
Publicando todas as 100 & 100 & 68 & 0,68 \\
\bottomrule
\end{tabular}
\end{center}
\vspace{-4pt}
{\small A faixa (intervalo de confiança de Wilson) mostra os valores compatíveis com o que medi.
Uma ressalva: a validação aninhada protege a escolha do limiar, mas os sinais e a forma da regra
foram escolhidos olhando os mesmos 100 rótulos. Por isso o 0,76 ainda pode ser otimista, e não
tenho como medir quanto.}

\bd{O que os rótulos mostraram sobre a convenção da equipe.}
\begin{itemize}
\item \bd{Temas gêmeos.} Oito temas aparecem duas vezes, com o mesmo título, em disciplinas
      diferentes (por exemplo, Lei de Execução Penal em Legislação Especial e em Processo Penal).
      Nas 6 questões da folha com tema gêmeo, a equipe usou o de Legislação Especial em 5; na Q0324
      marcou os dois. E não seguiu a disciplina informada na origem: a Q0091 veio como Direito
      Penal e foi para Lei de Execução Penal, em Legislação Especial. Na fase 1, eu mandava para
      revisão o gêmeo que não batia com a disciplina de origem; essa regra teria barrado um acerto.
      São só 6 casos.
\item \bd{Tema ou subtema.} A equipe rotula no tema em 53 das 92 questões e no subtema em 39, e
      indica temas extras em 39: questões que cobram mais de um assunto são comuns, e os extras dão
      mais de um jeito de acertar.
\item \bd{Fora da árvore:} 8 das 100 não têm tema adequado, e não encontrei nada em comum entre
      elas que permitisse detectá-las.
\item \bd{A régua pesa muito.} As mesmas 20 publicações acertam 0,95 pela régua do case, 0,90 se
      só o ramo contasse e 0,50 se fosse exigido o tema exato. O sistema acerta o assunto geral,
      não o detalhe.
\end{itemize}

% ================================= c) DIAGNOSTICO ==========================
\section{c) Diagnóstico}

\bd{O que atrapalha nos dados e nos temas.} Primeiro, as \bd{lixeiras}: temas genéricos que, na
classificação automática antiga, receberam centenas de questões que não eram deles. Quanto mais
cheio um tema está na origem, menos esse histórico informa. O próprio assunto de origem também
virou lixeira: ``Concurso de Crimes'' rotula 43 questões de temas variados. Segundo, os gêmeos,
que a régua não deixa trocar. Terceiro, a sobreposição entre uma lei especial e o instituto
correspondente do Código. Por último, um erro de catálogo: pela Lei 14.197/2021, os temas 5148 e
5156 são capítulos do tema 5155, não irmãos dele. Corrigi isso no entregável, sem trocar IDs.
Nenhuma publicação cai nesses temas, mas a mudança de texto mexeu em 18 decisões: 15 publicações
perdidas e 3 ganhas. Como a avaliação usa a árvore original, declaro o custo e levo a correção ao
professor.

\bd{Onde o sistema perde.} Nas 92 questões com tema: 20 publicadas certas e 1 publicada errada;
48 foram para revisão com o ramo certo já em primeiro lugar (18 passariam com um limiar mais baixo,
30 precisariam de uma confiança melhor); 18 tinham o ramo certo entre o 2º e o 10º lugar; e 5 nem
apareciam entre os dez primeiros. A busca já não é o gargalo principal: o que falta é o sistema
saber quando está certo.

\bd{Que questões o sistema ainda erra, e o que resolveria cada uma.}

\begin{center}
\small
\begin{tabular}{p{3.6cm}p{6.0cm}p{6.2cm}}
\toprule
Tipo de questão & O que acontece & O que resolveria \\
\midrule
Lei de Execução Penal & vai para Teoria da Pena (3 casos); o texto fala de remição ou falta
grave sem citar a lei & descrever o tema 4703 com o vocabulário da execução \\
Juizado Especial Criminal & vai para procedimento ou ação penal (3 casos), mesmo citando a Lei
9.099 & usar a lei citada como evidência, testada com rótulos novos \\
Gabarito em subtema & erra o ramo em 36\%, contra 19\% quando o gabarito é tema & mais rótulos
de subtema \\
Erro confiante & a Q0202 foi publicada errada com 0,97 de confiança: os quatro métodos leem o
mesmo texto e erram juntos & um sinal que não venha da busca, testado com rótulos novos \\
Fora da árvore & 8 casos sem nada em comum & revisão humana; criar os temas que faltam \\
\bottomrule
\end{tabular}
\end{center}

Legislação Especial concentra o problema: erra o ramo em 6 de 12 questões, contra 10\% em
Criminologia. Testei depois usar a lei citada como regra, mas a ideia veio de olhar esses erros e
tem poucos casos; ficou fora da entrega e é o primeiro teste da próxima rodada.

% ================================= d) PRODUCAO =============================
\section{d) Produção: 50 mil questões por semana}

\bd{Custo e tempo.} Classificar as 659 questões leva cerca de 40 segundos num processador comum;
50 mil levariam perto de uma hora, sem placa de vídeo e sem pagar API. O LLM só volta quando o
catálogo muda, para descrever os temas novos ou alterados, uma vez por tema. O custo que não escala
é a fila: com 17\% de publicação, uns 41 mil itens por semana iriam para revisão. Por isso a fila
deve ser ordenada pela confiança, com as questões mais perto do limiar primeiro: são as que mais
mudam de decisão com um rótulo novo.

\bd{Como eu colocaria em produção.} Com precisão estimada de 0,76, eu não deixaria o sistema
publicar sozinho no começo: uma em cada quatro publicadas estaria errada, que é o problema que a
plataforma quer resolver. Hoje ele serve como triagem: o ramo certo está em primeiro lugar em 74\%
das questões e entre os dez primeiros em 95\%, então o professor confirma uma sugestão ou escolhe
numa lista curta, em vez de procurar entre 375 temas. O caminho seria: (1) um período em que o
sistema só sugere e o professor decide; (2) publicação liberada por partes, só nos temas e faixas
de confiança em que a amostra auditada passar de 95\%; (3) auditoria separada das importações
novas, porque o sistema publica um pouco mais nas 300 questões recentes (18\%, 54 de 300) do que
nas antigas (15\%, 40 de 259), o que pode ser excesso de confiança.

\bd{Como a correção do professor vira dado de treino.} Cada revisão gera uma linha: questão, tema
sugerido, tema final e se a sugestão estava certa. Essa linha serve a quatro coisas:
\begin{enumerate}
\item é o rótulo de ``acertou ou errou'' que o modelo de confiança aprende, e esse é o gargalo
      atual;
\item amplia a folha-verdade, que recalibra o limiar e mede a precisão;
\item mostra quais temas acumulam correções, e esses ganham uma descrição melhor;
\item com alguns milhares de linhas, permite treinar um classificador direto sobre o texto.
\end{enumerate}

\needspace{7\baselineskip}
\bd{De quanto em quanto tempo o sistema se retreina.}
\begin{center}
\small
\begin{tabular}{lll}
\toprule
Peça & Quando & Condição para entrar em produção \\
\midrule
Modelo de confiança & toda semana, com as revisões novas & não piorar na validação aninhada \\
Limiar de publicação & a cada uns 100 rótulos novos & precisão aninhada de pelo menos 95\% \\
Descrição dos temas & quando o catálogo muda & passar na conferência automática \\
Classificador de texto & quando houver uns milhares de rótulos & superar a busca na validação \\
\bottomrule
\end{tabular}
\end{center}
\vspace{-3pt}
{\small Os volumes (100 rótulos, alguns milhares) são estimativas minhas, a confirmar com o uso.}

\bd{Como medir a qualidade.} Toda semana, três números: a precisão numa amostra de publicadas
conferida pelo professor, a fração de sugestões corrigidas na fila e os reportes de alunos por
tema. A amostra precisa ser grande. Com zero erros em $n$ questões conferidas, só dá para garantir
que o erro fica abaixo de cerca de $3/n$; para garantir 95\% de precisão, são precisas umas 60
publicações conferidas sem nenhum erro. Com 20, mesmo sem erro, só se garante 86\%.

\bd{O que eu faria com mais tempo.} Em ordem de retorno: (1) mais rótulos, começando pelas questões
que citam lei, pelas de Legislação Especial e pelas mais perto do limiar; (2) a regra da lei citada,
testada nesses rótulos novos; (3) o vocabulário da execução penal no tema 4703; (4) um sinal de
confiança que não venha da busca, já que hoje todos os sinais nascem da mesma comparação de texto
e erram juntos; (5) exemplos de questões por tema, propostos por um LLM e revisados por um
professor, para cobrir os 68 ramos que não aparecem na folha-verdade.

Uma conclusão atravessa o trabalho: o gargalo é o número de rótulos, não o modelo. Uma semana com
100 rótulos novos, bem escolhidos, renderia mais do que qualquer algoritmo novo.

% ================================= e) MEMO =================================
\clearpage
\section{e) Memo para o professor}
{\small\itshape Para você, especialista em Direito. Sem jargão de IA.}

\bd{O que eu construí}
\begin{itemize}
\item Um sistema que decide, para cada uma das 659 questões, se publica o tema ou manda para a sua
      revisão. Ele compara o texto da questão com a descrição dos temas e não consulta nenhum
      modelo de linguagem, como o ChatGPT, na hora de decidir. Publiquei 115 (17\%) e mandei 544
      para você.
\item Espero que umas 3 de cada 4 publicadas estejam certas. Não consigo prometer os 95\% pedidos:
      com 100 questões rotuladas, não há como provar isso.
\item Na prática, hoje o sistema é mais útil para sugerir o tema do que para publicar sozinho: o
      assunto certo aparece entre as dez primeiras opções em 95\% das questões.
\end{itemize}

\bd{O que encontrei}
\begin{itemize}
\item O sistema coloca o assunto certo em primeiro lugar em cerca de 74\% das questões; quando erra,
      em geral confunde temas vizinhos. Ele publica no tema principal, não no subtema, porque a
      régua aceita isso e é mais seguro.
\item A maior dificuldade é Legislação Especial: metade das questões vai para o tema
      correspondente do Código. Questões de execução da pena vão para Teoria da Pena, e questões de
      Juizado Especial vão para procedimento ou ação penal.
\item Quando a questão cita uma lei especial, isso costuma ser uma pista, mas ainda não usei como
      regra: são poucos casos e a ideia veio depois de eu ver os erros.
\item Nos temas com nome repetido em duas disciplinas, a sua equipe usou quase sempre o de
      Legislação Especial (5 de 6), e não a disciplina de onde a questão veio.
\end{itemize}

\bd{O que recomendo mudar nos temas}
\begin{itemize}
\item Dar nomes diferentes aos 8 pares de temas repetidos (por exemplo, citar a disciplina no
      título): 4684/6716, 5063/5466, 4688/6717, 4702/6723, 4704/6724, 4703/6734, 4705/6725 e
      4708/6728.
\item Deixar mais clara a fronteira entre Lei de Execução Penal e Teoria da Pena, e entre Juizado
      Especial Criminal e Procedimento Sumaríssimo.
\item Confirmar a correção que fiz em ``Instituições Democráticas'' (5148) e ``Funcionamento dos
      Serviços Essenciais'' (5156): pela Lei 14.197/2021, são capítulos de ``Estado Democrático de
      Direito'' (5155), não temas irmãos dele.
\item Rever os temas genéricos que hoje acumulam centenas de questões sem relação com eles.
\end{itemize}

\bd{O que preciso de você}
\begin{itemize}
\item Revisar primeiro as questões que eu ordenar na fila: as mais perto da decisão. Cada correção
      sua ensina o sistema e é o que permite, com o tempo, publicar mais com segurança.
\item Rotular de 200 a 300 questões, começando pelas que citam lei e pelas de Legislação Especial.
      É isso que permite, com o tempo, provar se o sistema chega a 95\%.
\item Confirmar qual dos dois nomes é o oficial em cada par de temas repetidos.
\end{itemize}

\end{document}
